\section{Experimental Setup}
\label{sec:experiments}

\subsection{Model and agent configuration}
All runs use Claude Opus 5.5 (\texttt{claude-opus-5-5}) through the Anthropic Messages API with adaptive
thinking and effort level \emph{high}, a soft budget of 12 tool calls, a hard cap of 30, at most 45 model turns
per task, and tool results truncated at 24{,}000 characters. Four tasks run in parallel. The local runtime
serves \bench through the same client interface as the competition runtime and scores each submission with
the scorer of \S\ref{sec:problem}. Like a server, it checks every call's arguments against the tool schema
and answers a non-conforming call with an error; the call still counts. This fixes what the \emph{no
validation} ablation measures: malformed calls cost runtime calls instead of being caught locally. The optimum $n^\star$ is set per task when the task is defined, as the
length of a minimal investigation path to the ground truth (2 to 5 calls, including the submission).

\subsection{Configurations}
The \emph{full} configuration enables every component of \S\ref{sec:method}. Five ablations each switch off
one component:
\begin{itemize}
\item \emph{no cache}: repeated calls are sent to the runtime again;
\item \emph{no validation}: arguments are not checked locally;
\item \emph{no grounding}: submissions citing unseen identifiers are sent without warning;
\item \emph{no budget}: no call-count feedback and no hard cap (the turn limit still applies);
\item \emph{no method}: the investigation-method prompt is replaced by a minimal prompt stating only the tool
rules (Appendix~\ref{app:prompts}).
\end{itemize}
Because LLM outputs vary between runs, each configuration is run three times on all 21 tasks, and we report
means with standard deviations across repeats. A single command runs the full configuration followed by the
five ablations on the same tasks. The first repeat was made with \bench v1.0 and rescored on v1.1
(\S\ref{sec:benchmark}); since the graphs are identical and scores depend only on the submission and the
number of calls, rescoring from the stored traces is exact. One run is excluded: in the first repeat of
\emph{no validation}, a malformed argument crashed the v1.0 runtime and aborted a task, which a schema-checking
runtime would have rejected instead; that configuration therefore has two repeats.

\subsection{Metrics}
Per task we record quality $Q$, efficiency $E$, reward $R$, tool calls spent, calls answered by the cache,
calls refused locally, wall-clock time and tokens (input, output, cache reads). Aggregates per configuration:
mean $Q$, $E$, $R$ and calls; verdict accuracy; the false-positive rate (benign tasks judged malicious) and the
false-negative rate (malicious tasks judged benign); each also broken down by operating system and by expected
verdict. All figures in \S\ref{sec:results} are regenerated from the run folders by a script in the release.

\subsection{Pilot}
Before the Linux tasks existed, we ran the full configuration once on the \ResPilotTasks{} Windows tasks (then
\bench v0). The pilot used the same model and prompt and is reported as scored at the time, with the v1.0
ground truth. The revision to v1.1 leaves the Windows malicious ground truth unchanged; the three Windows benign
tasks gained small neutral sets (one or two lineage ids each), so the pilot's benign quality can only be
understated.

\subsection{Implementation and testing}
The agent is written in Python and installs into the competition's public reference client. Its test suite
(25 tests, in addition to the reference client's own) covers the catalog normalisation, every safeguard and
ablation switch, the field aliasing, parallel execution against a simulated runtime with a task limit
(including through the reference client's own transport layer), deadline handling, the graph builder on
Windows and Linux events, scoring and the ground-truth rules, the runtime's rejection of malformed calls, the
benchmark pack round trip, offline rescoring and a scripted run on the real recordings.
